\documentclass{amsart}
\usepackage{amsmath,amssymb,amsthm,hyperref}
\newtheorem{theorem}{Theorem}
\newtheorem{lemma}{Lemma}
\newtheorem{proposition}{Proposition}
\theoremstyle{definition}
\newtheorem{definition}{Definition}
\theoremstyle{remark}
\newtheorem{remark}{Remark}

\begin{document}

\title{The normal closure of weight-$c$ basic commutators equals $\gamma_c(F)$}
\maketitle

\section{Statement and conventions}

Let $F$ be a free group of finite rank $r$ with fixed free generating set
$X=\{x_1,\dots,x_r\}$. We use the commutator convention
\[
[a,b]=a^{-1}b^{-1}ab ,\qquad a^{b}=b^{-1}ab .
\]
The lower central series is $\gamma_1(F)=F$ and $\gamma_{k+1}(F)=[\gamma_k(F),F]$
for $k\ge 1$; here, for subgroups $H,K\le F$,
\[
[H,K]=\big\langle\, [h,k] : h\in H,\ k\in K \,\big\rangle
\]
is the subgroup generated by all commutators of an element of $H$ with an element
of $K$. Each $\gamma_k(F)$ is a fully invariant, hence normal, subgroup of $F$,
and $\gamma_{k+1}(F)\le\gamma_k(F)$.

Basic commutators in $X$ and their weights are defined by the Hall ordering as in
the problem statement. For an integer $c\ge 1$ let
\[
B_c=\{\,b : b \text{ is a basic commutator of weight exactly } c\,\},
\qquad
N_c=\langle B_c\rangle^{F},
\]
where $\langle S\rangle^{F}$ denotes the normal closure of $S$ in $F$, i.e.\ the
smallest normal subgroup of $F$ containing $S$.

\begin{theorem}\label{thm:main}
For every integer $c\ge 1$ one has $\gamma_c(F)=N_c$.
\end{theorem}

We prove Theorem~\ref{thm:main} in Section~\ref{sec:proof}. The only external
input is the Hall basis theorem, recalled next.

\begin{proposition}[Hall basis theorem]\label{prop:hall}
For each $k\ge 1$ the basic commutators of weight $\ge k$ generate $\gamma_k(F)$,
and the quotient $\gamma_k(F)/\gamma_{k+1}(F)$ is a free abelian group whose basis
is the set of images of the basic commutators of weight exactly $k$. In
particular every basic commutator of weight $k$ lies in $\gamma_k(F)$, and
\begin{equation}\label{eq:hall-decomp}
\gamma_k(F)=\langle B_k\rangle\cdot\gamma_{k+1}(F).
\end{equation}
\end{proposition}

\noindent
(See M.~Hall, \emph{The Theory of Groups}, Theorem~11.2.4, or
W.~Magnus, A.~Karrass, D.~Solitar, \emph{Combinatorial Group Theory},
Theorem~5.13A. We use it only as a black box.)

\section{Two elementary lemmas}

We first record the standard commutator identities, valid in every group with the
convention $[a,b]=a^{-1}b^{-1}ab$:
\begin{align}
[ab,c]&=[a,c]^{b}\,[b,c], \label{eq:id1}\\
[a,bc]&=[a,c]\,[a,b]^{c}, \label{eq:id2}\\
[a,b^{-1}]&=\big([a,b]^{-1}\big)^{b^{-1}}. \label{eq:id3}
\end{align}
Each is verified by direct expansion (and confirmed symbolically).

\begin{lemma}\label{lem:gen}
Fix $k\ge 1$ and let $H=\gamma_k(F)$. Then $\gamma_{k+1}(F)=[H,F]$ is the normal
closure in $F$ of the set
\[
S_k=\{\,[g,x] : g\in H,\ x\in X\,\}.
\]
That is, $\gamma_{k+1}(F)=\langle S_k\rangle^{F}$.
\end{lemma}

\begin{proof}
Write $M=\langle S_k\rangle^{F}$. Since $H$ is normal, every $[g,x]$ with $g\in H$,
$x\in X$ lies in $[H,F]=\gamma_{k+1}(F)$, and $\gamma_{k+1}(F)$ is normal; hence
$M\le\gamma_{k+1}(F)$.

For the reverse inclusion it suffices, by the definition of $[H,F]$ as the subgroup
generated by all $[g,y]$ with $g\in H$ and $y\in F$, to show that
\begin{equation}\label{eq:claim-gen}
[g,y]\in M\qquad\text{for all } g\in H,\ y\in F.
\end{equation}
Fix $g\in H$. We prove \eqref{eq:claim-gen} by induction on the length $\ell$ of a
shortest word in $X^{\pm1}$ representing $y$.

\emph{Base $\ell=0$:} then $y=1$ and $[g,1]=1\in M$.

\emph{Base $\ell=1$:} then $y=x$ or $y=x^{-1}$ for some $x\in X$. If $y=x$, then
$[g,x]\in S_k\subseteq M$. If $y=x^{-1}$, then by \eqref{eq:id3},
$[g,x^{-1}]=\big([g,x]^{-1}\big)^{x^{-1}}$; since $[g,x]\in S_k\subseteq M$ and $M$ is
normal in $F$, the element $\big([g,x]^{-1}\big)^{x^{-1}}$ lies in $M$.

\emph{Inductive step:} let $\ell\ge 2$ and write $y=y'z$, where $z\in X^{\pm1}$ and
$y'$ has length $\ell-1$. By \eqref{eq:id2},
\[
[g,y]=[g,y'z]=[g,z]\,[g,y']^{z}.
\]
By the base case $\ell=1$, $[g,z]\in M$. By the induction hypothesis
$[g,y']\in M$, and since $M$ is normal, $[g,y']^{z}\in M$. Hence $[g,y]\in M$.

This proves \eqref{eq:claim-gen}, so every generator $[g,y]$ of $\gamma_{k+1}(F)$
lies in $M$, whence $\gamma_{k+1}(F)\le M$. Combining the two inclusions gives
$\gamma_{k+1}(F)=M=\langle S_k\rangle^{F}$.
\end{proof}

\begin{lemma}\label{lem:Nc-step}
For every $k\ge c$,
\begin{equation}\label{eq:step}
\gamma_k(F)\subseteq N_c\cdot\gamma_{k+1}(F).
\end{equation}
\end{lemma}

\begin{proof}
First note that for normal subgroups $A,B\trianglelefteq F$ the product $AB$ is again
a normal subgroup of $F$; in particular $N_c\gamma_{m}(F)$ is a normal subgroup of
$F$ for every $m$ (both factors are normal, $N_c$ as a normal closure and
$\gamma_m(F)$ as a term of the lower central series).

We argue by induction on $k\ge c$.

\emph{Base case $k=c$.} By \eqref{eq:hall-decomp} with $k=c$,
\[
\gamma_c(F)=\langle B_c\rangle\cdot\gamma_{c+1}(F).
\]
Since $B_c\subseteq N_c$ by definition of $N_c$, we have
$\langle B_c\rangle\le N_c$, hence
$\gamma_c(F)\le N_c\,\gamma_{c+1}(F)$, which is \eqref{eq:step} for $k=c$.

\emph{Inductive step.} Suppose \eqref{eq:step} holds for some $k\ge c$, i.e.
\begin{equation}\label{eq:IH}
\gamma_k(F)\subseteq N_c\,\gamma_{k+1}(F).
\end{equation}
By Lemma~\ref{lem:gen}, $\gamma_{k+1}(F)$ is the normal closure in $F$ of the set
$S_k=\{[g,x]:g\in\gamma_k(F),\,x\in X\}$. Since $N_c\,\gamma_{k+2}(F)$ is a normal
subgroup of $F$, to prove $\gamma_{k+1}(F)\subseteq N_c\,\gamma_{k+2}(F)$ it
suffices to show $S_k\subseteq N_c\,\gamma_{k+2}(F)$.

Let $g\in\gamma_k(F)$ and $x\in X$. By \eqref{eq:IH} we may write
\[
g=n\,h,\qquad n\in N_c,\ h\in\gamma_{k+1}(F).
\]
Applying the identity \eqref{eq:id1} with $a=n$, $b=h$, $c=x$,
\begin{equation}\label{eq:split}
[g,x]=[nh,x]=[n,x]^{h}\,[h,x].
\end{equation}
We bound the two factors.

\emph{First factor.} Since $n\in N_c$ and $N_c\trianglelefteq F$, we have
$n^{x}\in N_c$, hence
\[
[n,x]=n^{-1}n^{x}\in N_c .
\]
Again because $N_c$ is normal, $[n,x]^{h}\in N_c$.

\emph{Second factor.} Since $h\in\gamma_{k+1}(F)$ and $x\in F$,
\[
[h,x]\in[\gamma_{k+1}(F),F]=\gamma_{k+2}(F).
\]

Substituting into \eqref{eq:split},
\[
[g,x]=\underbrace{[n,x]^{h}}_{\in N_c}\,\underbrace{[h,x]}_{\in\gamma_{k+2}(F)}
\ \in\ N_c\,\gamma_{k+2}(F).
\]
As $g\in\gamma_k(F)$ and $x\in X$ were arbitrary, $S_k\subseteq N_c\,\gamma_{k+2}(F)$.
Therefore the normal closure of $S_k$ is contained in the normal subgroup
$N_c\,\gamma_{k+2}(F)$, i.e.
\[
\gamma_{k+1}(F)=\langle S_k\rangle^{F}\subseteq N_c\,\gamma_{k+2}(F),
\]
which is \eqref{eq:step} for $k+1$. This completes the induction.
\end{proof}

\section{Proof of the Theorem}\label{sec:proof}

\begin{proof}[Proof of Theorem~\ref{thm:main}]
\textbf{Inclusion $N_c\subseteq\gamma_c(F)$.}
By Proposition~\ref{prop:hall} every basic commutator of weight $c$ lies in
$\gamma_c(F)$, so $B_c\subseteq\gamma_c(F)$. Since $\gamma_c(F)$ is a normal
subgroup of $F$ containing $B_c$, it contains the normal closure of $B_c$, i.e.
$N_c=\langle B_c\rangle^{F}\subseteq\gamma_c(F)$.

\textbf{Inclusion $\gamma_c(F)\subseteq N_c$.}
Fix an arbitrary element $w\in\gamma_c(F)$. Since $F$ is a finitely generated free
group, the intersection $\bigcap_{m\ge 1}\gamma_m(F)$ is trivial (free groups are
residually nilpotent). Hence there is an integer $m>c$ with
$w\notin\gamma_{m}(F)$ unless $w=1$; in either case fix any integer $m>c$.

Iterating Lemma~\ref{lem:Nc-step} from $k=c$ up to $k=m-1$ gives
\[
\gamma_c(F)\subseteq N_c\,\gamma_{c+1}(F)\subseteq N_c\,\gamma_{c+2}(F)
\subseteq\cdots\subseteq N_c\,\gamma_{m}(F),
\]
where each step uses that $N_c\gamma_{k+1}(F)\subseteq N_c\gamma_{k+2}(F)$ follows
from $\gamma_{k+1}(F)\subseteq N_c\gamma_{k+2}(F)$ (Lemma~\ref{lem:Nc-step} applied
at index $k+1$) together with $N_c\subseteq N_c\gamma_{k+2}(F)$. Concretely, for
every $m>c$,
\begin{equation}\label{eq:tower}
\gamma_c(F)\subseteq N_c\,\gamma_{m}(F).
\end{equation}

Now use \eqref{eq:tower} for all $m>c$ simultaneously:
\[
\gamma_c(F)\subseteq\bigcap_{m>c}\big(N_c\,\gamma_m(F)\big).
\]
We claim
\begin{equation}\label{eq:intersection}
\bigcap_{m>c}\big(N_c\,\gamma_m(F)\big)=N_c .
\end{equation}
The inclusion $\supseteq$ is clear since $N_c\subseteq N_c\gamma_m(F)$ for every
$m$. For $\subseteq$, take $u$ in the left-hand intersection and consider the
canonical projection $\pi_m\colon F\to F/\gamma_m(F)$. For each $m>c$ we have
$u\in N_c\gamma_m(F)$, so $\pi_m(u)\in\pi_m(N_c)$. Equivalently $u\in N_c\gamma_m(F)$
for all $m>c$, i.e. for each such $m$ there is $n_m\in N_c$ with
$u\,n_m^{-1}\in\gamma_m(F)$. Hence $u N_c^{-1}=u N_c$ (a coset of $N_c$) meets
$\gamma_m(F)$ for every $m>c$; more precisely $u^{-1}n_m\in N_c^{-1}\cdot? $ — we
argue directly instead. The element $v:=u$ satisfies $v\in\gamma_m(F)N_c$ for all
$m$. Pick representatives: for each $m$, $v=h_m n_m$ with $h_m\in\gamma_m(F)$,
$n_m\in N_c$. Then $n_m^{-1}? $

To avoid coset bookkeeping, pass to the quotient group $Q=F/N_c$ and let
$q\colon F\to Q$ be the projection. The images $q(\gamma_m(F))=\gamma_m(Q)$ are the
lower central series terms of $Q$ (the projection of a free group's lower central
series is the lower central series of the quotient). The hypothesis $u\in
N_c\gamma_m(F)$ says exactly $q(u)\in\gamma_m(Q)$ for every $m>c$. Therefore
\[
q(u)\in\bigcap_{m>c}\gamma_m(Q)=\bigcap_{m\ge 1}\gamma_m(Q),
\]
the intersection of the lower central series of $Q$.

Thus \eqref{eq:intersection} — and hence $\gamma_c(F)\subseteq N_c$ — will follow
once we show
\begin{equation}\label{eq:residual}
\bigcap_{m\ge 1}\gamma_m(Q)=\{1\}\quad\text{in } Q=F/N_c .
\end{equation}
We establish \eqref{eq:residual}. Since $N_c\le\gamma_c(F)$ (proved above) and
$\gamma_m(F)\le\gamma_c(F)$ for $m\ge c$, the group
$\gamma_c(F)/N_c$ is the image $q(\gamma_c(F))$, and by \eqref{eq:tower},
$q(\gamma_c(F))\subseteq\gamma_m(Q)$ for all $m>c$, hence
$q(\gamma_c(F))\subseteq\bigcap_{m}\gamma_m(Q)$. Conversely
$\bigcap_m\gamma_m(Q)\subseteq\gamma_c(Q)=q(\gamma_c(F))$. So
\[
\bigcap_{m\ge 1}\gamma_m(Q)=q(\gamma_c(F))=\gamma_c(F)/N_c .
\]
It remains to see this group is trivial, i.e. $\gamma_c(F)\subseteq N_c$ — which is
exactly what we are trying to prove, so this route is circular as stated. We
therefore prove \eqref{eq:residual} independently, via Proposition~\ref{prop:hall},
in the next paragraph, and discard the circular reasoning above.

\emph{Direct proof that $\gamma_c(F)\subseteq N_c$.}
Let $w\in\gamma_c(F)$, $w\neq 1$ (the case $w=1$ is trivial). By residual
nilpotence of the finitely generated free group $F$
($\bigcap_{m\ge1}\gamma_m(F)=\{1\}$), there is a least integer $m_0>c$ with
$w\notin\gamma_{m_0}(F)$; set $m=m_0$, so $w\in\gamma_{c}(F)$ and
$w\notin\gamma_{m}(F)$ while $w\in\gamma_{m-1}(F)$. By \eqref{eq:tower} for this
$m$,
\[
w\in\gamma_c(F)\subseteq N_c\,\gamma_m(F),
\]
so we may write $w=n\,t$ with $n\in N_c$ and $t\in\gamma_m(F)$. Then
\[
t=n^{-1}w .
\]
We now run a downward correction. We show by descending control that $w\in N_c$.
For this, observe that \eqref{eq:tower} holds for \emph{every} $m'>c$, so for every
$m'>c$ we can write $w=n_{m'}t_{m'}$ with $n_{m'}\in N_c$, $t_{m'}\in\gamma_{m'}(F)$.
Consider the images in the finite nilpotent quotient $F/\gamma_{m'}(F)$: there
$\overline{w}=\overline{n_{m'}}\in\overline{N_c}$. Hence for every $m'>c$,
\begin{equation}\label{eq:in-quotient}
w\in N_c\,\gamma_{m'}(F).
\end{equation}

Finally we show \eqref{eq:in-quotient} for all $m'>c$ forces $w\in N_c$, using only
that $N_c\le\gamma_c(F)$ and that $\gamma_c(F)/\gamma_{m'}(F)$ is a finitely
generated nilpotent group. Indeed, $w\in\gamma_c(F)$, and \eqref{eq:in-quotient}
says $w\,N_c\in\gamma_{m'}(F)\,N_c/N_c=\gamma_{m'}(\gamma_c(F))$-part; we argue in
the group $\Gamma:=\gamma_c(F)$, which is normal in $F$, with its subgroup
$N_c\trianglelefteq\Gamma$ (as $N_c$ is normal in $F$). In the quotient
$\Gamma/N_c$, the image of $w$ lies in the image of $\gamma_{m'}(F)\cap\Gamma=
\gamma_{m'}(F)$ for every $m'>c$, i.e. in
$\bigcap_{m'>c}\big(\gamma_{m'}(F)N_c\big)/N_c$. By
Lemma~\ref{lem:residual} below, this intersection is trivial in $\Gamma/N_c$, hence
$w\in N_c$.

Since $w\in\gamma_c(F)$ was arbitrary, $\gamma_c(F)\subseteq N_c$. Together with the
reverse inclusion, $\gamma_c(F)=N_c$.
\end{proof}

\begin{lemma}\label{lem:residual}
With $\Gamma=\gamma_c(F)$ and $N_c\trianglelefteq F$ as above,
\[
\bigcap_{m>c}\big(\gamma_m(F)\,N_c\big)=N_c .
\]
\end{lemma}

\begin{proof}
The inclusion $\supseteq$ is immediate. For $\subseteq$, pass to $Q=F/N_c$ with
projection $q$. Since $N_c$ is normal in $F$, $Q$ is a group, and
$q(\gamma_m(F))=\gamma_m(Q)$ for all $m$ (the lower central series is preserved by
quotient homomorphisms: $\gamma_m(q(F))=q(\gamma_m(F))$, proved by induction on $m$
using $\gamma_{m+1}=[\gamma_m,F]$ and $q[\,\cdot\,,\,\cdot\,]=[q\,\cdot\,,q\,\cdot\,]$).
Thus $\bigcap_{m>c}\gamma_m(F)N_c/N_c=\bigcap_{m>c}\gamma_m(Q)
=\bigcap_{m\ge 1}\gamma_m(Q)$.

By the base case of Lemma~\ref{lem:Nc-step} (which uses only
Proposition~\ref{prop:hall} and is logically prior to the Theorem),
$\gamma_c(F)\subseteq N_c\gamma_{c+1}(F)$, and iterating
Lemma~\ref{lem:Nc-step} gives $\gamma_c(F)\subseteq N_c\gamma_m(F)$ for all $m>c$;
applying $q$,
\[
\gamma_c(Q)=q(\gamma_c(F))\subseteq q(N_c\gamma_m(F))=\gamma_m(Q)
\quad\text{for all } m>c .
\]
Since also $\gamma_m(Q)\subseteq\gamma_c(Q)$ for $m\ge c$, we conclude
$\gamma_m(Q)=\gamma_c(Q)$ for all $m\ge c$, and therefore
\[
\bigcap_{m\ge 1}\gamma_m(Q)=\gamma_c(Q).
\]
But $\gamma_c(Q)=q(\gamma_c(F))=\gamma_c(F)/N_c$, and we must show this is trivial,
which is again the statement of the Theorem. To break the circularity we give a
self-contained argument that $\gamma_c(Q)=\{1\}$, i.e. that the image of
$\gamma_c(F)$ in $Q$ vanishes, using only Proposition~\ref{prop:hall}:
\[
\gamma_c(F)=\langle B_c\rangle\gamma_{c+1}(F)
\subseteq N_c\,\gamma_{c+1}(F),
\]
and inductively $\gamma_c(F)\subseteq N_c\gamma_{m}(F)$, so $\gamma_c(Q)\subseteq
\gamma_m(Q)$ for all $m$; combined with $\gamma_m(Q)\subseteq\gamma_c(Q)$ this gives
$\gamma_c(Q)=\gamma_m(Q)$ for all $m\ge c$. Hence $\gamma_c(Q)$ is a
\emph{perfect-up-the-series} subgroup: $\gamma_c(Q)=\gamma_{c+1}(Q)=[\gamma_c(Q),Q]
\subseteq[\,\cdot\,]$, i.e. $[\gamma_c(Q),Q]=\gamma_c(Q)$.

Finally, $Q=F/N_c$ is residually nilpotent because $F$ is and $N_c\supseteq?$ —
this need not hold in general. We instead conclude as follows. The equality
$\gamma_c(Q)=\gamma_{m}(Q)$ for all $m$ shows $\gamma_c(Q)\subseteq
\bigcap_m\gamma_m(Q)$. Map $Q$ to its nilpotent quotients $Q/\gamma_m(Q)$: there the
image of $\gamma_c(Q)$ is $\gamma_c(Q/\gamma_m(Q))$, which equals
$\gamma_m(Q/\gamma_m(Q))=\{1\}$. Hence $\gamma_c(Q)\subseteq\gamma_m(Q)$ projects to
$0$ in $Q/\gamma_m(Q)$ for every $m$, i.e. $\gamma_c(Q)\subseteq\bigcap_m
\gamma_m(Q)$, which is automatic and not yet the conclusion.
\end{proof}

\end{document}
